\section{\texorpdfstring{Sčítání a odčítání v jiných soustavách}{Scitani a odcitani v jinych soustavach}}

% =====================================================================================================================
	\begin{frame}
    \frametitle{Sčítání}

		\begin{minipage}[t][\textheight][t]{0.49\textwidth}
      Sčítání v soustavě o základu 10
      \vspace{10mm}
      
			\boldmath
			\begin{tabular}{ p{20mm} c c c}
				1. člen $X$  & 1 & 2 & 2 \\
				2. člen $Y$  & 2 & 9 & 1 \\
        přenos    & \textcolor{red}{1} &   &   \\
        \hline
        výsledek  & 4 & 1 & 3 \\
			\end{tabular}
      
      \vspace{10mm}
      \begin{itemize}
        \item Přenos při $x + y > 10$
        \item Zapisujeme $x + y - 10$
      \end{itemize}
      
		\end{minipage}
		\begin{minipage}[t][\textheight][t]{0.5\textwidth}
      Sčítání v soustavě o základu 8
      \vspace{10mm}
      
			\boldmath
			\begin{tabular}{ p{20mm} c c c}
				1. člen $X$  & 1 & 7 & 2 \\
				2. člen $Y$  & 4 & 4 & 3 \\
        přenos    & \textcolor{red}{1} &   &   \\
        \hline
        výsledek  & 6 & 3 & 5 \\
			\end{tabular}
      
      \vspace{10mm}
      \begin{itemize}
        \item Přenos při $x + y > 8$
        \item Zapisujeme $x + y - 8$
      \end{itemize}
      
		\end{minipage}
		
	\end{frame}

% =====================================================================================================================
	\begin{frame}
    \frametitle{Sčítání vlastnosti}
		\boldmath
    
    \begin{itemize}
      \item Ve všech soustavách stejný postup.
      \item Při sčítání dvou sčítanců může být přenos maximálně 1.
      \item Při sčítání dvou sčítanců může být výsledek o jednu pozici širší než jsou oba sčítance.
    \end{itemize}
    
    Později bude dávat smysl rozšíření obou sčítanců nulami vlevo na možnou délku výsledku.
    
    \vspace{4mm}
    \textbf{Příklad pro binární hodnoty}:
    \vspace{2mm}
      
    \boldmath
    \begin{center}
      \begin{tabular}{ p{20mm} c c c c c c c}
        1. člen $X$  & \textcolor{blue}{0} & 1 & 1 & 1 & 0 & 1 & 0 \\
        2. člen $Y$  & \textcolor{blue}{0} & 1 & 0 & 1 & 0 & 1 & 1 \\
        přenos       & \textcolor{red}{1}  & \textcolor{red}{1} & \textcolor{red}{1} &   & \textcolor{red}{1} &   &   \\
        \hline
        výsledek     &                  1  & 1 & 0 & 0 & 1 & 0 & 1 \\
      \end{tabular}
    \end{center}
    
	\end{frame}
  
% =====================================================================================================================
	\begin{frame}
    \frametitle{Odčítání}

		\begin{minipage}[t][\textheight][t]{0.49\textwidth}
      Odčítání v soustavě o základu 10
      \vspace{10mm}
      
			\boldmath
			\begin{tabular}{ p{20mm} c c c c}
				1. člen $X$  &   & 2 & 9 & 1 \\
				2. člen $Y$  & - & 1 & 2 & 2 \\
        výpůjčka     &   &  & \textcolor{red}{-1} &   \\
        \hline
        výsledek     &   & 1 & 6 & 9 \\
			\end{tabular}
      
      \vspace{10mm}
      \begin{itemize}
        \item Přenos při $x + y > 10$
        \item Zapisujeme $x - y + 10$
      \end{itemize}
      
		\end{minipage}
		\begin{minipage}[t][\textheight][t]{0.5\textwidth}
      Odčítání v soustavě o základu 8
      \vspace{10mm}
      
			\boldmath
			\begin{tabular}{ p{20mm} c c c c}
				1. člen $X$  &   & 4 & 4 & 3 \\
				2. člen $Y$  & - & 1 & 7 & 2 \\
        výpůjčka     &   & \textcolor{red}{-1} &   &   \\
        \hline
        výsledek     &   & 2 & 5 & 1 \\
			\end{tabular}
      
      \vspace{10mm}
      \begin{itemize}
        \item Výpůjčka při $x - y < 0$
        \item Zapisujeme $x - y + 8$
      \end{itemize}
      
		\end{minipage}
		
	\end{frame}
  
% =====================================================================================================================
	\begin{frame}
    \frametitle{Odčítání vlastnosti}
		\boldmath
    
    \begin{itemize}
      \item Ve všech soustavách stejný postup.
      \item Při rozdílu dvou členů může být výpůjčka maximálně 1.
      \item Při přímém vyjádření znaménka musí být menšenec větší než menšitel.
      \item Pro větší menšitel než menšenec stále odečítáme menší číslo od většího a měníme znaménko u výsledku.
      \item Výsledek může být jen kratší, viz příklad.
    \end{itemize}
    
    \vspace{4mm}
    \textbf{Příklad pro binární hodnoty}:
    \vspace{2mm}
      
    \boldmath
    \begin{center}
      \begin{tabular}{ p{20mm} c c c c c c c c}
        1. člen $X$  &   & 1 & 1 & 1 & 0 & 1 & 0 \\
        2. člen $Y$  & - & 1 & 0 & 1 & 0 & 1 & 1 \\
        výpůjčka     &   &   & \textcolor{red}{-1}  & \textcolor{red}{-1} & \textcolor{red}{-1}  & \textcolor{red}{-1}  &   \\
        \hline
        výsledek     &   & 0 & 0 & 0 & 1 & 1 & 1 \\
      \end{tabular}
    \end{center}
    
	\end{frame}